\section{Experiment}\label{sec:experiment}
\subsection{Experiment Setup}\label{sub:exp_setup}
%然后的话，那个实验那一章，要从第十页开始写，写到第十四页中间才可以，因为最后那一章结论那一章只能占半页，

%(5.1, 0) Experiment Setup
% Firstly, we summarize the experiment setup, including the baselines we compare with, the benchmarks we evaluate on, the implementation details of our method and baselines, and the evaluation metrics we use.

\subsubsection{Baselines}
We implemented several novel multi-modal memory methods as baselines for comparison:

\begin{enumerate}
    \item \textbf{M3Agent~\cite{long2025M3Agent}: } A novel multimodal agent framework handling with long-term memory.
    \item \textbf{EgoBulter~\cite{yang2025EgoLife}: } Integration of an omni-modal model trained on egocentric datasets (EgoGPT) and a retrieval-based component that supports answering ultra-long-context questions (EgoRAG).
    \item \textbf{Ego-R1\cite{tian2025EgoR1}: } The reinforcement learning enhanced agent that can learn to select the most relevant video clips for question answering.
    \item \textbf{HippoMM~\cite{lin2025HippoMM}: } A biologically-inspired framework that applies hippocampal mechanisms for multimodal understanding.
    % \item \textbf{MemVerse~\cite{liu2025MemVerse}: } An open-source solution designed to provide continuous and multimodal memory for AI agents.
    \item \textbf{Open Source VLMs: } We also compare with several open-source VLMs, including Qwen Series~\cite{yang2024Qwen2,qwen2025Qwen25,xu2025qwen25omni, yang2025Qwen3}, Intern Series~\cite{zhu2025InternVL3, wang2025InternVideo25} and LLaVA Series~\cite{liu2023LLaVA, li2024LLaVAOneVision, zhang2025LLaVAVideo}, which are the most popular open-source VLMs with strong performance.
\end{enumerate}

\subsubsection{Benchmarks}

The experiment is executed on the following three benchmarks:

\begin{enumerate}
    \item \textbf{Egoschema~\cite{mangalam2023EgoSchema}. } This benchmark contains 5031 questions on 5031 egocentric video clips from Ego4D, which is the largest egocentric video question-answering benchmark so far. 500 hundred answers are released.
    \item \textbf{EgoLifeQA~\cite{yang2025EgoLife}. } A behchmark proposed along with the EgoLife dataset, which contains 500 questions on the egocentric video from one person's perspective in EgoLife.
    \item \textbf{Ego-R1 Bench~\cite{tian2025EgoR1}. } A benchmark proposed along with the Ego-R1 method, which contains 300 questions on the EgoLife dataset (50 question for each of the 6 videos in the EgoLife dataset).
\end{enumerate}

% 	(5.1.3, 1, 4) Evaluation Mode
\subsubsection{Evaluation Mode}
As these benchmarks are multiple-choice questions, each of the questions $\mathcal{Q}$ is composed of a question stem $q^0$ and several options $\{o_i\}$, and the answer is the index of the correct option $a$ (so the content of the correct option is $o_a$).
Existing works only evaluate the three benchmarks above in the multiple-choice mode, i.e., they input the question stem and all the options into the models, and let the models give its choice as the answer.
Such evaluation have three main flaws: (1) it is not realistic, as in real life the user would not provide options along with their question;
(2) the options contain information about the question that is not included in the question stem, so the model might be hinted by the options and may not fully understand video to answer;
(3) when model realizes that this is a multiple-choice question, even it has no understanding of the video and question, it can just guess one option and will still reach an accuracy of about 20\% to 25\%.

Based on these consideration, we evaluate the methods by open-ended question-answering mode besides the traditional multiple-choice mode. (1) For multiple-choice mode, where we input the question stem and all the options into the model, and let the model print its choice as the answer, $r^c$;
\begin{equation}
    q^c = concat(q^0, o^1, o^2...) \quad r^c := \mathcal{R}(\mathbf{M}_\tau, q^c)
\end{equation}
(2) For open-ended mode, where we only input the question stem into the models, and let the model generate an answer in natural language $r^e$.
\begin{equation}
    q^e = q^0 \quad r^e := \mathcal{R}(\mathbf{M}_\tau, q^e)
\end{equation}
% In daily practice, the user would not provide options along with their question, so the open-ended mode is more realistic, but only few works evaluate in this way.
% We provide both evaluation results in these two modes for a more comprehensive comparison.

\subsubsection{Metrics}

There are mainly two types of evaluation metrics in this experiment: accuracy metrics and efficiency metrics.

% (5.1.4, 1) Accuracy
% 	(5.1.4, 1, 1) Overall
The accuracy metrics indicate the quality of the answer generated by the model, and we have the following six accuracy metrics:
%   (5.1.4, 1, 2) List

\textbf{(1) $\alpha^c$ Option-question Accuracy.} For multiple-choice mode, we directly check if the option chosen by the model is the correct answer.
\begin{equation}
    \alpha^c := (r^c == a)
\end{equation}

\textbf{(2) $\alpha^e_E$ Closest Option from Encoding.} For open-ended question answering mode, we first encode the model-generated answer and all the options into vectors, and then check if the option closest to the model-generated answer in vector space is the correct answer.
\begin{equation}
    \alpha^e_E := (argmax_i ||\mathcal{S}(r^e, o_i)|| == a)
\end{equation}

\textbf{(3) $\theta^e$ Similarity to Ground Truth.} We also directly check the similarity between the model-generated answer and the ground truth answer in the vector space, to see how close the model-generated answer is to the ground truth answer.
\begin{equation}
    \theta^e := \mathcal{S}(r^e, o_a)
\end{equation}

\textbf{(4) $\alpha^e_J$ LLM Judgement.} Besides the direct check in vector space, we also use a large language model to read the question and both the model-generated answer and the ground truth, and judge whether the two answers are consistent.
\begin{equation}
    \alpha^e_J := (LLM(r^e, o_a, \text{`tell if same'}) == \text{`same'})
\end{equation}
M3Agent\cite{long2025M3Agent} uses $\alpha^e_J$ to evaluate the response accuracy.

\textbf{(5) $\alpha^e_C$ LLM Choosing.} We also use a large language model to read the model-generated answer and all the options, and choose which option is closest to the model-generated answer, and check if this option is the correct answer.
\begin{equation}
    \alpha^e_C := (LLM(r^e, \{o_i\}, \text{`select closest'}) == a)
\end{equation}

\textbf{(6) $\alpha^{e+}_C$ LLM Choosing with Abstention.} The last metric is a stricter version of the above metric, where we allow the large language model to abstain from choosing any option if it thinks none of the options is close enough to the model-generated answer, and check if the large language model's choice is correct.
\begin{equation}
    \alpha^{e+}_C := (LLM(r^e, \{o_i\} \cup \{None\}, \text{`select closest'}) == a)
\end{equation}




Among these metrics, the first accuracy metric is for multiple-choice mode, and the rest five accuracy metrics are for open-ended question-answering mode.
During evaluation, we choose Qwen3 Embedding (0.6B) as the encoder for the second and third accuracy metrics, and we choose Qwen3 (235B) as the judge for the fourth and selector for the fifth and sixth accuracy metrics.

Efficiency metrics evaluate the time and space consumption of the method. We calculate following six efficiency metrics:
\begin{enumerate}
    \item \textbf{$\delta^c$ Multiple-choice Delay}: the inference time for multiple-choice mode.
    \item \textbf{$\delta^e$ Open-ended Delay}: the inference time for open-ended mode.
    \item \textbf{$\delta^c_\tau=\delta^c/\tau$ Multiple-choice Time Rate}: ratio of inference time for multiple choice mode to reference video duration
    \item \textbf{$\delta^e_\tau=\delta^e/\tau$ Open-ended Time Rate}: ratio of inference time for open-ended mode to reference video duration
    \item \textbf{$\delta^m_\tau=\delta^m/\tau$ Memorize Time Rate}: ratio of memory time to reference video duration
    \item \textbf{$\eta^{\mathbf{M}}_{\mathbf{V}}=\eta(\mathbf{M})/\eta(\mathbf{V})$ Compression Rate}: ratio of memory file size to reference video size
\end{enumerate}
Delay, denoted as $\delta$, is the average of the two mode's delay. Delay rate, denoted as $\delta_\tau$, is the average of the two mode's delay rate.
After calculating these metrics on each question-answer pair, we will average the metrics across all QA pairs to get the final accuracy metrics for each method on each benchmark.


\begin{table}[htbp]
    \centering
    \fontsize{7}{10}\selectfont
    % \scriptsize
    \setlength{\tabcolsep}{6pt} % 调整列间距
    \renewcommand{\arraystretch}{1.4} % 调整行高
    \centering
    \begin{tabular}{r|p{0.3cm}<{\centering}p{0.3cm}<{\centering}p{0.3cm}<{\centering}p{0.35cm}<{\centering}|p{0.3cm}<{\centering}p{0.3cm}<{\centering}p{0.3cm}<{\centering}p{0.35cm}<{\centering}|p{0.3cm}<{\centering}p{0.3cm}<{\centering}p{0.3cm}<{\centering}p{0.35cm}<{\centering}}
        \textbf{Benchmark} &
            \multicolumn{4}{c|}{\fontsize{9}{12}\selectfont Egoschema} &
            \multicolumn{4}{c|}{\fontsize{9}{12}\selectfont EgoLifeQA} &
            \multicolumn{4}{c}{\fontsize{9}{12}\selectfont EgoR1Bench}\\ \hline \hline
        \textbf{Metric} 
        & $\overline{\alpha^c}$ & $\overline{\alpha^e_E}$ & $\overline{\theta^e}$ & $\overline{\alpha^e_J}$
        & $\overline{\alpha^c}$ & $\overline{\alpha^e_E}$ & $\overline{\theta^e}$ & $\overline{\alpha^e_J}$
        & $\overline{\alpha^c}$ & $\overline{\alpha^e_E}$ & $\overline{\theta^e}$ & $\overline{\alpha^e_J}$
        \\
        Unit 
        & \% & \% & -- & \%
        & \% & \% & -- & \%
        & \% & \% & -- & \%
        \\  \hline
        %\textit{Open Source Models} & & & & & & & & & & \\
        Qwen3VL(32B)~\cite{yang2025Qwen3}               & 68.4                     & 59.2                     & .608                     & 51.8\cellcolor{second!60}    & 36.4\cellcolor{third!60}  & 23.8                     & .312                     &  29.8\cellcolor{second!60}& 33.7                     & 23.3                     & .325                     & 28.0\cellcolor{second!60}\\
        Qwen2.5VL(72B)~\cite{qwen2025Qwen25}            & 72.8                     & 59.8\cellcolor{third!60} & .608                     & 44.4\cellcolor{third!60}     & 29.6                      & 24.8                     & .275                     &  9.60                     & 34.7                     & 24.0                     & .287                     & 21.0                     \\
        LLaVA-Video(72B)~\cite{zhang2025LLaVAVideo}     & 74.0                     & 55.0                     & .689\cellcolor{third!60} & 38.0                         & 35.4                      & 21.2                     & .599\cellcolor{second!60}&  6.80                     & 41.0\cellcolor{third!60} & 19.7                     & .578\cellcolor{first!60} & 17.3                     \\
        LLaVA-OV(72B)~\cite{li2024LLaVAOneVision}       & 66.0                     & 49.2                     & .673                     & 37.2                         & 38.0\cellcolor{second!60} & 25.2                     & .606\cellcolor{first!60} &  8.80                     & 41.7\cellcolor{second!60}& 26.0\cellcolor{third!60} & .578\cellcolor{second!60}& 18.7                     \\
        InternVL3.5(38B)~\cite{zhu2025InternVL3}        & 80.2\cellcolor{second!60}& 59.8\cellcolor{second!60}& .719\cellcolor{first!60} & 40.6                         & 35.8                      & 24.0                     & .464                     &  6.23                     & 35.7                     & 24.0                     & .503                     & 17.7                     \\
        InternVideo2.5(8B)~\cite{wang2025InternVideo25} & 74.4\cellcolor{third!60} & 53.0                     & .640                     & 31.4                         & 28.0                      & 23.8                     & .478                     &  5.60                     & 30.0                     & 22.0                     & .502                     & 14.0                     \\
        EgoGPT(7B)~\cite{yang2025EgoLife}               & 57.6                     & 48.6                     & .660                     & 33.8                         & 35.0                      & 26.0\cellcolor{third!60} & .594\cellcolor{third!60} &  9.00                     & 30.0                     & 22.0                     & .570                     & 21.7                     \\
        LongVA(7B)~\cite{zhang2024LongVA}               & 49.2                     & 43.0                     & .597                     & 34.2                         & 34.2                      & 24.2                     & .338                     &  3.60                     & 30.7                     & 24.0                     & .370                     & 14.3                     \\
        \hline
        %\textit{Open Source Methods} & & & & & & & & & & & & & \\
        EgoButler(7B)~\cite{yang2025EgoLife}            & 38.5                     & 46.2                     & .525                     & 23.1                         & 27.7                      & 26.4                     & .415                     &  4.40                     & 21.0                     & 22.1                     & .467                     & 13.3                     \\
        Ego-R1(3B)$^{\dag}$~\cite{tian2025EgoR1}        & 65.4                     & 49.2                     & .567                     & 23.4                         & 35.0                      & 24.0                     & .450                     &  5.20                     & 45.0                     & 24.6                     & .394                     & 12.0                     \\
        M3Agent(70B)~\cite{long2025M3Agent}             & 59.4                     & 50.0                     & .594                     & 31.4                         & 31.4                      & 22.2                     & .451                     &  22.6\cellcolor{third!60} & 31.7                     & 30.7\cellcolor{second!60}& .460                     & 21.7\cellcolor{third!60} \\
        HippoMM(*)$^{\dag}$~\cite{lin2025HippoMM}       & 55.6                     & 25.0                     & .628                     & 37.4                         & 16.7                      & 26.1\cellcolor{second!60}& .305                     &  4.40                     & 14.3                     & 9.33                     & .347                     & 5.00                     \\
        % MemVerse~\cite{liu2025MemVerse} & & & & & & & & & & \\
        EgoVidMem(*)                                    & 88.6\cellcolor{first!60} & 74.2\cellcolor{first!60} & .692\cellcolor{second!60}& 56.2\cellcolor{first!60}     & 55.0\cellcolor{first!60}  & 34.2\cellcolor{first!60} & .451                     &  30.4\cellcolor{first!60} & 59.0\cellcolor{first!60} & 38.7\cellcolor{first!60} & .571\cellcolor{third!60} & 29.7\cellcolor{first!60} \\ \hline \hline
    \end{tabular}
    \caption{
        Experimental results of several methods on three egocentric benchmarks.
        %解释一下单次MLLM的使用方法
        Eight MLLMs (above the line) and four memorizing framework (below the line) are compared with ours.
        These MLLMs take each ultra-long video with 64 uniformly sampled frames as input, and answer the questions in a zero-shot manner.
        The $\dag$ symbol after Ego-R1 and HippoMM indicates that these methods require the original video as input.
        We list the size of each model, and the asterisk symbol following HippoMM and EgoVidMem indicates these two methods do not directly provide own finetuned models, but uses other large models as modules.
        The best, second best and third best results in each column are highlighted with pink, orange and light yellow.
    }
    \label{tab:exp_acc}
\end{table} 


\begin{table}[htbp]
    \centering
    \fontsize{7}{10}\selectfont
    % \scriptsize
    \setlength{\tabcolsep}{6pt}
    \renewcommand{\arraystretch}{1.4}
    \begin{tabular}{r|p{0.4cm}<{\centering}p{0.4cm}<{\centering}p{0.4cm}<{\centering}p{0.4cm}<{\centering}|p{0.4cm}<{\centering}p{0.4cm}<{\centering}p{0.4cm}<{\centering}p{0.4cm}<{\centering}|p{0.4cm}<{\centering}p{0.4cm}<{\centering}p{0.4cm}<{\centering}p{0.4cm}<{\centering}}
        \textbf{Benchmark} &
            \multicolumn{4}{c|}{\fontsize{9}{12}\selectfont Egoschema} &
            \multicolumn{4}{c|}{\fontsize{9}{12}\selectfont EgoLifeQA} &
            \multicolumn{4}{c}{\fontsize{9}{12}\selectfont EgoR1Bench}\\ \hline \hline
        \textbf{Metric} 
        & $\overline{\delta}$ & $\overline{\delta_\tau}$ & $\overline{\delta^m_\tau}$ & $\overline{\eta^{\mathcal{M}}_{\mathcal{V}}}$
        & $\overline{\delta}$ & $\overline{\delta_\tau}$ & $\overline{\delta^m_\tau}$ & $\overline{\eta^{\mathcal{M}}_{\mathcal{V}}}$
        & $\overline{\delta}$ & $\overline{\delta_\tau}$ & $\overline{\delta^m_\tau}$ & $\overline{\eta^{\mathcal{M}}_{\mathcal{V}}}$
        \\
        Unit 
        & \multicolumn{3}{c}{s$\quad \times10^{-1} \quad $ --} & --
        & \multicolumn{3}{c}{s$\quad \times10^{-4} \quad $ --} & --
        & \multicolumn{3}{c}{s$\quad \times10^{-4} \quad $ --} & --
        \\
        \hline
        EgoButler(7B)~\cite{yang2025EgoLife}
        & 2.16\cellcolor{first!60} & .120\cellcolor{first!60} & .333\cellcolor{first!60}  & .021\cellcolor{first!60}
        & 1.71\cellcolor{first!60} & .672\cellcolor{first!60} & .192\cellcolor{second!60} & .001\cellcolor{first!60}
        & 2.06\cellcolor{first!60} & .161\cellcolor{first!60} & .170\cellcolor{second!60} & .001\cellcolor{first!60}
        \\
        Ego-R1(3B)$^{\dag}$~\cite{tian2025EgoR1}
        & 21.6 & 1.31 & --- & 1.04
        & 33.8\cellcolor{third!60} & 16.0\cellcolor{third!60} & --- & 1.02
        & 28.6\cellcolor{third!60} & 16.1\cellcolor{third!60} & --- & 1.02
        \\
        M3Agent(70B)~\cite{long2025M3Agent}
        & 34.9 & 1.93 & .510 & .135\cellcolor{third!60}
        & 86.9 & 34.2 & 2.55 & .017\cellcolor{second!60}
        & 79.6 & 56.0 & 2.56 & .015\cellcolor{second!60}
        \\
        HippoMM(*)$^{\dag}$~\cite{lin2025HippoMM}
        & 14.4\cellcolor{third!60} & .800\cellcolor{third!60} & .396\cellcolor{third!60} & 1.85
        & 134. & 52.8 & .049\cellcolor{first!60} & 1.23
        & 85.2 & 33.6 & .049\cellcolor{first!60} & 1.25
        \\
        EgoVidMem(*)
        & 6.03\cellcolor{second!60} & .335\cellcolor{second!60} & .362\cellcolor{second!60} & .048\cellcolor{second!60}
        & 16.3\cellcolor{second!60} & 7.73\cellcolor{second!60} & .381\cellcolor{third!60} & .070\cellcolor{third!60}
        & 12.8\cellcolor{second!60} & 5.36\cellcolor{second!60} & .341\cellcolor{third!60} & .063\cellcolor{third!60}
        \\
        \hline
    \end{tabular}
    \caption{
        Efficiency performance of five memorizing frameworks on three egocentric benchmarks.
        % Same as Table~\ref{tab:exp_acc}, the three best values are shaded.
        The $\dag$ symbol indicates that these method requires the original video as input, as a result, their compression rate $\overline{\eta^{\mathcal{M}}_{\mathcal{V}}}$ are larger than 1.0.
        The hierachical rag construction process of Ego-R1 is not provided, so we cannot calculate $\overline{\delta^m_\tau}$ for it.
        Considering the long duration of the original video, we also shift the unit of $\overline{\delta_\tau}$ by $\times10^{-1}$ for EgoSchema and $\times10^{-4}$ for the others to make the values more readable.
    }
    \label{tab:exp_per}
\end{table}


\subsection{Experimental Result}\label{sub:exp_result}


Table~\ref{tab:exp_acc} shows the accuracy results of different methods on the three benchmarks. Due to the space limit, only the first four metrics are listed in the table, and the last two metrics (which are less significant) will be shown in the appendix.

%解读一下表格？
%准确率领先
Through the numbers in Table~\ref{tab:exp_acc}, \textbf{EgoVidMem} performs better than other baselines in most of the accuracy metrics.
%Egoschema比其他两个好，因为它短。
Our method performs best on the Egoschema benchmark.
For the other two benchmarks which have much longer videos and more complex questions, our method strongly improves the baselines by 10-20\%.
%虽然我们的方法在问答题评价范式之下相对较好，但不出所料，问答题回答范式的准确率有大幅下降。这对“这些方法是否真正理解了问题和视频内容“提出了不小的质疑。
Although our method performs relatively well in the open ended mode, it is not surprising that all the methods' accuracy drop significantly compared to the multiple-choice mode. This raises a fundamental doubt of "whether these methods truly understand the question and video content".
%在LLM评价者模式下，准确率甚至降到了个位数，这说明这些方法
For the LLM Judgement $\alpha^e_J$, some accuracy is even dropped to a single digit, which indicates that these methods' responses are not really similar to the ground truth.
%similarity指标的显著性相较于另外两项比较差，其有效性还有待进一步验证。
The similarity metric $\theta^e$ is less consistent with the other two open-ended metrics, so its validity might require further verification. 


\begin{table}[htbp]
    \centering
    \fontsize{7}{10}\selectfont
    \setlength{\tabcolsep}{6pt} % 调整列间距
    \renewcommand{\arraystretch}{1.4} % 调整行高
    \centering
    \begin{tabular}{r|p{0.5cm}<{\centering}p{0.5cm}<{\centering}p{0.5cm}<{\centering}|p{0.5cm}<{\centering}p{0.5cm}<{\centering}p{0.5cm}<{\centering}|p{0.5cm}<{\centering}p{0.5cm}<{\centering}p{0.5cm}<{\centering}}
        \textbf{Benchmark} &
            \multicolumn{3}{c|}{\fontsize{10}{12}\selectfont Egoschema} &
            \multicolumn{3}{c|}{\fontsize{10}{12}\selectfont EgoLifeQA} &
            \multicolumn{3}{c}{\fontsize{10}{12}\selectfont EgoR1Bench}\\ \hline \hline
        \textbf{Metric}
        & $\overline{\alpha^c}$ & $\overline{\alpha^e_E}$ & $\overline{\eta^{\mathcal{M}}_{\mathcal{V}}}$
        & $\overline{\alpha^c}$ & $\overline{\alpha^e_E}$ & $\overline{\eta^{\mathcal{M}}_{\mathcal{V}}}$
        & $\overline{\alpha^c}$ & $\overline{\alpha^e_E}$ & $\overline{\eta^{\mathcal{M}}_{\mathcal{V}}}$
        \\
        Unit 
        & \% & \% & --
        & \% & \% & --
        & \% & \% & --
        \\  \hline
        $\mu$=8, k=6     &                          &                          &                          &                          &                          &                          &                          &                          &                          \\
        EgoVidMem        & 88.6                     & 74.2                     & .050                     & 55.0                     & 34.2                     & .070                     & 59.0                     & 38.7                     & .063                     \\ \hline
        w/o Images       & -0.6                     & -0.4                     & -.05\cellcolor{first!60} & -8.2                     & -3.6                     & -.07\cellcolor{first!60} & -9.7                     & -14.\cellcolor{second!60}& -.06\cellcolor{first!60} \\ 
        w/o Summary      & -60.\cellcolor{first!60} & -55.\cellcolor{first!60} & -0.0                     & -8.6                     & -5.4\cellcolor{third!60} & -0.0                     & -10.                     & -10.                     & -0.0                     \\ 
        w/o Transcripts  & -0.0                     & -0.0                     & -0.0                     & -27.\cellcolor{first!60} & -12.\cellcolor{first!60} & -0.0                     & -28.\cellcolor{first!60} & -18.\cellcolor{first!60} & -0.0                     \\ \hline
        w/o Context      & -8.2\cellcolor{second!60}& -4.4                     & -.01                     & -9.4\cellcolor{second!60}& -0.4                     & -0.0                     & -19.\cellcolor{second!60}& -8.7                     & -0.0                     \\ \hline
        k=4              & -1.8                     & -6.2\cellcolor{third!60} & -0.0                     & -2.6                     & -4.0                     & -0.0                     & -8.7                     & -8.3                     & -0.0                     \\ 
        k=2              & -1.6                     & -7.6\cellcolor{second!60}& -0.0                     & -5.0                     & -3.4                     & -0.0                     & -11.                     & -11.                     & -0.0                     \\ \hline
        $\mu$=8          & -0.0                     & -0.0                     & -0.0                     & -8.2                     & -2.2                     & -0.0                     & -2.7                     & -10.                     & -0.0                     \\
        $\mu$=12         & -0.0                     & -0.0                     & -0.0                     & -8.8\cellcolor{third!60} & -8.6\cellcolor{second!60}& -0.0                     & -12.\cellcolor{third!60} & -13.\cellcolor{third!60} & -0.0                     \\ \hline
        Text Encoder     & -2.6\cellcolor{third!60} & +0.8                     & -.05\cellcolor{first!60} & -8.4                     & -3.6                     & -.07\cellcolor{first!60} & -9.2                     & -12.                     & -.06\cellcolor{first!60} \\ \hline
    \end{tabular}
    \caption{
        Ablation studies on the EgoVidMem's hyper parameters and its memory unit components. 
        The results are reported as the change in performance compared to the first row, with color shaded to highlight the three most significant drops in performance, similar to Table~\ref{tab:exp_acc}.
        Besides the clarified modification in the first column, all the other settings are consistent with the basic version of EgoVidMem.
        The basic version of EgoVidMem (first row) is evaluated with $\mu$=8 and k=6. Its encoder is Qwen3 VL Embedding 8B, while the last row 'text encoder' utilized Qwen3 8B.
    }
    \label{tab:ablation}
\end{table}


Besides the accuracy, Table~\ref{tab:exp_per} shows the efficiency results of different memorization frameworks. These MLLMs in Table~\ref{tab:exp_acc} do not contain memorizing process and they don't have memorization time and compression rate, thus, we do not compare with them in this part.
EgoBulter\cite{yang2025EgoLife} proposed from the paper of EgoLife takes the least space and time among all these. However, its accuracy is not comparable to other baselines.
For M3Agent\cite{long2025M3Agent}, it takes much longer time than other method and even longer than the reference video duration ($\overline{\delta^m_\Delta} > 1$), which is not acceptable in streaming input scenario.
M3Agent\cite{long2025M3Agent}, HippoMM\cite{lin2025HippoMM} and Ego-R1\cite{tian2025EgoR1} utilize chain of thought process to answer questions, which is time-consuming and might not be practical for application.
Generally, our method takes practical time and space for memorization and question-answering, whose accuracy is also better than other baselines.

\subsection{Ablation Study}\label{sub:exp_ablation}


We evaluate the following variants of our method as ablation study.
(1) We ablate the three components of our memory.
(2) We generate the memory without providing context information, to see how the context information affects the performance.
(3) We decrease the $k$ value in the RAG process. 
(4) We increase the transcript encoding granularity $\mu$.
(5) We change the encoder $\mathcal{E}$ in the Responder to a textual encoder, Qwen3 Embedding (8B).
This version's method also can't see the images in the memory.

In Table~\ref{tab:ablation}, we list three significant metrics: multiple-choice accuracy ($\overline{\alpha^c}$), open-ended accuracy ($\overline{\alpha^e_E}$), and compression rate ($\overline{\eta^{\mathcal{M}}_{\mathcal{V}}}$) of these ablation variants. The full results of all the metrics will be shown in the appendix.
Egoschema doesn't contain transcripts, so the summary contains almost all the information. Thus, the accuracy drops significantly when we ablate the summary.
However, for the other two benchmarks, the transcript plays a more essential role.
Besides, the encoding granularity and context information also have significant effect on the performance.
Furthermore, the compression rate of our method is around 5\% to 7\%, which is acceptable for long-term memory in streaming input scenario.
%在其中，图片占据了几乎全部的空间，比总结和字幕的空间大两个数量级。
Among the three components of our memory module, the image component occupies almost all the space, which is two orders of magnitude larger than the summary and transcript components.
